% GENERATED FILE. Edit canonical lessons, then run npm run generate:books.
% canonical-source-hash: fnv1a64:e1e1635f6d6da1ca
% canonical-lessons: RU-C342-obsuzhdat, RU-C342-obsuzhdenie, RU-C342-diskussiya, RU-C342-vyskazyvatsya, RU-C342-delitsya

\chapter{Talking it over}
\label{ch:ru-talking-it-over}
\hlchaptermodality{342}

\begin{chapteropening}
\textbf{By the end of this chapter:} \emph{I can discuss a question, take part in a discussion or a debate, give my opinion and share my ideas.}

Complete the last lesson of chapter 342: I can discuss a question, take part in a discussion or a debate, give my opinion and share my ideas.
\end{chapteropening}

\section[\texorpdfstring{\ru{обсуждать} (obsuzhdát)}{обсуждать (obsuzhdát)}]{\ru{обсуждать} (obsuzhdát) — to discuss}

\label{lesson:RU-C342-obsuzhdat}



\begin{quote}
Before the new one: say the Russian for to insist, then the Russian for to admit.
\end{quote}

\subsection*{You'll want to know: \ru{обсуждать}}
\textbf{\ru{обсуждать}} (\emph{obsuzhdát}) — \textquotedblleft{}to discuss\textquotedblright{}.

\begin{quote}
\ru{Вчера} \ru{мы} \ru{обсуждали} \ru{план} \ru{на} \ru{лето}.
\end{quote}

\begin{quote}
\emph{Vcherá my obsuzhdáli plan na léto.}
\end{quote}

\begin{quote}
Yesterday we discussed the plan for the summer.
\end{quote}

Imperfective; perfective \emph{\ru{обсудить}}. It takes the accusative with no preposition: \emph{\ru{обсуждать} \ru{фильм}}.

\begin{cousinweb}
\textbf{\ru{обсуждать}} is built from \emph{\ru{об}-} \textquotedblleft{}around\textquotedblright{} and \emph{\ru{судить}} \textquotedblleft{}to judge\textquotedblright{}, on \emph{\ru{суд}} \textquotedblleft{}a court, a judgement\textquotedblright{}, inherited from Proto-Slavic \emph{*sǫdŭ} (EtymDB). To discuss is to judge a thing from all sides.
\end{cousinweb}

\begin{grammarlens}[title={What you've built}]
One more word for telling, explaining and describing at B1.
\end{grammarlens}

\subsection*{Your turn}
\begin{itemize}
\raggedright
  \item \emph{Say it:} \emph{obsuzhdát}
  \item \emph{Say it:} \emph{obsuzhdát}, once more
  \item \emph{Say it:} a sentence of your own with \emph{obsuzhdát}
  \item \emph{Recall:} say the Russian for to insist, then the Russian for to admit, then say \emph{obsuzhdát} again
\end{itemize}

\begin{tcolorbox}[breakable,colback=teal!4,colframe=teal!35!black,title={Before you move on}]
What does \ru{обсуждать} mean? (\textquotedblleft{}To discuss\textquotedblright{}.) Say it once more.
\end{tcolorbox}
\section[\texorpdfstring{\ru{обсуждение} (obsuzhdéniye)}{обсуждение (obsuzhdéniye)}]{\ru{обсуждение} (obsuzhdéniye) — a discussion}

\label{lesson:RU-C342-obsuzhdenie}



\begin{quote}
Before the new one: say the Russian for to admit, then the Russian for to discuss.
\end{quote}

\subsection*{You'll want to know: \ru{обсуждение}}
\textbf{\ru{обсуждение}} (\emph{obsuzhdéniye}) — \textquotedblleft{}a discussion\textquotedblright{}.

\begin{quote}
\ru{Обсуждение} \ru{продолжалось} \ru{два} \ru{часа}.
\end{quote}

\begin{quote}
\emph{Obsuzhdéniye prodolzhálos dva chasá.}
\end{quote}

\begin{quote}
The discussion went on for two hours.
\end{quote}

Neuter. \emph{\ru{Вынести} \ru{на} \ru{обсуждение}} — to put up for discussion.

\begin{cousinweb}
\textbf{\ru{обсуждение}} is built from \emph{\ru{обсуждать}} \textquotedblleft{}to discuss\textquotedblright{} and the suffix \emph{-\ru{ение}} (EtymDB); the root is \emph{\ru{суд}}, Proto-Slavic \emph{*sǫdŭ} (EtymDB).
\end{cousinweb}

\begin{grammarlens}[title={What you've built}]
One more word for telling, explaining and describing at B1.
\end{grammarlens}

\subsection*{Your turn}
\begin{itemize}
\raggedright
  \item \emph{Say it:} \emph{obsuzhdéniye}
  \item \emph{Say it:} \emph{obsuzhdéniye}, once more
  \item \emph{Say it:} a sentence of your own with \emph{obsuzhdéniye}
  \item \emph{Recall:} say the Russian for to admit, then the Russian for to discuss, then say \emph{obsuzhdéniye} again
\end{itemize}

\begin{tcolorbox}[breakable,colback=teal!4,colframe=teal!35!black,title={Before you move on}]
What does \ru{обсуждение} mean? (\textquotedblleft{}A discussion\textquotedblright{}.) Say it once more.
\end{tcolorbox}
\section[\texorpdfstring{\ru{дискуссия} (diskússiya)}{дискуссия (diskússiya)}]{\ru{дискуссия} (diskússiya) — a debate, an organised discussion}

\label{lesson:RU-C342-diskussiya}



\begin{quote}
Before the new one: say the Russian for to discuss, then the Russian for a discussion.
\end{quote}

\subsection*{You'll want to know: \ru{дискуссия}}
\textbf{\ru{дискуссия}} (\emph{diskússiya}) — \textquotedblleft{}a debate, an organised discussion\textquotedblright{}.

\begin{quote}
\ru{После} \ru{фильма} \ru{началась} \ru{интересная} \ru{дискуссия}.
\end{quote}

\begin{quote}
\emph{Pósle fílma nachalás interésnaya diskússiya.}
\end{quote}

\begin{quote}
After the film an interesting debate began.
\end{quote}

Feminine. More formal than \emph{\ru{обсуждение}}: a public or organised debate.

\begin{cousinweb}
\textbf{\ru{дискуссия}} was borrowed from Latin \emph{discussiō} (EtymDB) \textquotedblleft{}a shaking apart, an examination\textquotedblright{}, from \emph{discutere} \textquotedblleft{}to strike apart, to examine\textquotedblright{} (Lewis \& Short).
\end{cousinweb}

\begin{grammarlens}[title={What you've built}]
One more word for telling, explaining and describing at B1.
\end{grammarlens}

\subsection*{Your turn}
\begin{itemize}
\raggedright
  \item \emph{Say it:} \emph{diskússiya}
  \item \emph{Say it:} \emph{diskússiya}, once more
  \item \emph{Say it:} a sentence of your own with \emph{diskússiya}
  \item \emph{Recall:} say the Russian for to discuss, then the Russian for a discussion, then say \emph{diskússiya} again
\end{itemize}

\begin{tcolorbox}[breakable,colback=teal!4,colframe=teal!35!black,title={Before you move on}]
What does \ru{дискуссия} mean? (\textquotedblleft{}A debate, an organised discussion\textquotedblright{}.) Say it once more.
\end{tcolorbox}
\section[\texorpdfstring{\ru{высказываться} (vyskázyvatsya)}{высказываться (vyskázyvatsya)}]{\ru{высказываться} (vyskázyvatsya) — to speak one’s mind, to give one’s opinion}

\label{lesson:RU-C342-vyskazyvatsya}



\begin{quote}
Before the new one: say the Russian for a discussion, then the Russian for a debate.
\end{quote}

\subsection*{You'll want to know: \ru{высказываться}}
\textbf{\ru{высказываться}} (\emph{vyskázyvatsya}) — \textquotedblleft{}to speak one’s mind, to give one’s opinion\textquotedblright{}.

\begin{quote}
\ru{На} \ru{собрании} \ru{каждый} \ru{мог} \ru{высказаться}.
\end{quote}

\begin{quote}
\emph{Na sobránii kázhdyy mog výskazatsya.}
\end{quote}

\begin{quote}
At the meeting everyone could give their opinion.
\end{quote}

Imperfective; perfective \emph{\ru{высказаться}} (stressed on the first syllable). \emph{\ru{Высказываться} \ru{за} / \ru{против}} — to speak for or against.

\begin{cousinweb}
\textbf{\ru{высказываться}} is built from \emph{\ru{вы}-} \textquotedblleft{}out\textquotedblright{} and \emph{\ru{сказать}} \textquotedblleft{}to say\textquotedblright{}; the root is \emph{\ru{казать}} \textquotedblleft{}to show\textquotedblright{}, inherited from Proto-Slavic \emph{*kazati} (EtymDB). To speak your mind is to say it out.
\end{cousinweb}

\begin{grammarlens}[title={What you've built}]
One more word for telling, explaining and describing at B1.
\end{grammarlens}

\subsection*{Your turn}
\begin{itemize}
\raggedright
  \item \emph{Say it:} \emph{vyskázyvatsya}
  \item \emph{Say it:} \emph{vyskázyvatsya}, once more
  \item \emph{Say it:} a sentence of your own with \emph{vyskázyvatsya}
  \item \emph{Recall:} say the Russian for a discussion, then the Russian for a debate, then say \emph{vyskázyvatsya} again
\end{itemize}

\begin{tcolorbox}[breakable,colback=teal!4,colframe=teal!35!black,title={Before you move on}]
What does \ru{высказываться} mean? (\textquotedblleft{}To speak one’s mind, to give one’s opinion\textquotedblright{}.) Say it once more.
\end{tcolorbox}
\section[\texorpdfstring{\ru{делиться} (delítsya)}{делиться (delítsya)}]{\ru{делиться} (delítsya) — to share}

\label{lesson:RU-C342-delitsya}



\begin{quote}
Before the new one: say the Russian for a debate, then the Russian for to speak one’s mind.
\end{quote}

\subsection*{You'll want to know: \ru{делиться}}
\textbf{\ru{делиться}} (\emph{delítsya}) — \textquotedblleft{}to share\textquotedblright{}.

\begin{quote}
\ru{Она} \ru{делится} \ru{с} \ru{нами} \ru{своими} \ru{идеями}.
\end{quote}

\begin{quote}
\emph{Oná délitsya s námi svoími idéyami.}
\end{quote}

\begin{quote}
She shares her ideas with us.
\end{quote}

Imperfective; present \emph{\ru{делюсь}, \ru{делишься}}; perfective \emph{\ru{поделиться}}. \emph{\ru{Делиться} \ru{чем} \ru{с} \ru{кем}}: the thing shared goes in the instrumental.

\begin{cousinweb}
\textbf{\ru{делиться}} is the reflexive of \emph{\ru{делить}} \textquotedblleft{}to divide, to share\textquotedblright{}, inherited from Proto-Slavic \emph{*děliti} \textquotedblleft{}to divide\textquotedblright{}, built on \emph{*dělŭ} \textquotedblleft{}a part\textquotedblright{} (EtymDB); EtymDB sets Proto-Germanic \emph{*dailą} \textquotedblleft{}part, deal\textquotedblright{} beside it.
\end{cousinweb}

\begin{grammarlens}[title={What you've built}]
One more word for telling, explaining and describing at B1.
\end{grammarlens}

\subsection*{Your turn}
\begin{itemize}
\raggedright
  \item \emph{Say it:} \emph{delítsya}
  \item \emph{Say it:} \emph{delítsya}, once more
  \item \emph{Say it:} a sentence of your own with \emph{delítsya}
  \item \emph{Recall:} say the Russian for a debate, then the Russian for to speak one’s mind, then say \emph{delítsya} again
\end{itemize}

\begin{tcolorbox}[breakable,colback=teal!4,colframe=teal!35!black,title={Before you move on}]
What does \ru{делиться} mean? (\textquotedblleft{}To share\textquotedblright{}.) Say it once more.
\end{tcolorbox}
